\chapter{Algebraic Foundations}
\label{chapter:Algebraic-Foundations}

%&& a &= b & \text{Proof \ref{proof:}} \label{equation:}

bla bla bla operations past numbers which are addition-like or multiplication like

\section{Group-Like Structures}

Groups generalize binary operations with addition- or multiplication-like properties.

\begin{definition}[Monoid]
	A monoid $\{ M, \cdot \}$ is a set $M$ with a binary operation $\cdot: G \times G \rightarrow G$ which satisfies the following properties.
	\begin{enumerate}
		\item Associativity
		\begin{equation}
			\forall m, n, p \in M: m \cdot (n \cdot p) = (m \cdot n) \cdot p
		\end{equation}
		\item Identity element
		\begin{equation}
			\exists i \in M: \forall m \in M: i \cdot m = m \cdot i = m
		\end{equation}
	\end{enumerate}
\end{definition}

\begin{definition}[Group]
	A group $\{ G, \cdot \}$ is a monoid which also has inverse elements.
	\begin{equation}
		\forall a \in G: \exists b \in A: a \cdot b = i
	\end{equation}
\end{definition}

\begin{definition}[Abelian Group]
	An Abelian group $\{A, \cdot\}$ is a group which also has commutativity
	\begin{equation}
		\forall a, b \in A: a \cdot b = b \cdot a
	\end{equation}
\end{definition}

In the same spirit of generalizing addition and multiplication, the following two notation conventions are often used for such operations.

\begin{itemize}
	\item Additive: An addition-like operation, almost always an Abelian group, with $+$ for the operation, $0$ for the identity, and $-a$ for any inverses. For a group, subtraction can be defined.
	\begin{equation}
		a - b := a + (-b)
	\end{equation}
	\item Multiplicative: A multiplication-like operation, less often an Abelian group, with $\times$, $\cdot$, or $*$ for the operation, $1$ for the identity, and $a^{-1}$ for any inverses. For a group, division can be defined.
	\begin{equation}
		a / b := a \cdot b^{-1}
	\end{equation}
\end{itemize}

\subsection{Group Actions}

Actions generalize operations between arbitrary sets and monoids. For example, displacements are an action on positions. Compositions of displacements form an additive Abelian group using ``tip to tail'' addition. If one composition of displacements equals another composition, then they should get you to the same final destination.

\begin{definition}[Left Action]
	A left action $\circ: M \times X \rightarrow X$ is a binary operation between a monoid $\{ M, \cdot \}$ with identity $i$ and a set $X$ which satisfies the following properties.
	\begin{enumerate}
		\item Left identity
		\begin{equation}
			\forall x \in X: i \circ x = x
		\end{equation}
		\item Left compatibility
		\begin{equation}
			\forall m, n \in M, x \in X: m \circ (n \circ x) = (m \cdot n) \circ x
		\end{equation}
	\end{enumerate}
\end{definition}

\begin{definition}[Right Action]
	A right action $\circ: X \times M \rightarrow X$ is a binary operation between a set $X$ and a monoid $\{ M, \cdot \}$ with identity $i$ which satisfies the following properties.
	\begin{enumerate}
		\item Right identity
		\begin{equation}
			\forall x \in X: x \circ i = x
		\end{equation}
		\item Right compatibility
		\begin{equation}
			\forall x \in X, m, n \in G: (x \circ m) \circ n = x \circ (m \cdot n)
		\end{equation}
	\end{enumerate}
\end{definition}

\section{Ring-Like Structures}

Rings generalize cases where there is both an addition- and multiplication-like operation. Multiplication may or may not be invertible, but if it is, cannot inlcude the additive identity 0. See Proof \ref{proof:div_by_zero}.

\begin{definition}[Ring]
	A ring $\{ R, +, \cdot \}$ is a set $R$ with an additive Abelian group $\{ R, + \}$ and a multiplicative monoid $\{ R, \cdot \}$ which satisfy the following properties.
	\begin{enumerate}
		\item Left distributivity
		\begin{equation}
			\forall r, s, t \in R: r \cdot (s + t) = r \cdot s + r \cdot t
		\end{equation}
		\item Right distributivity
		\begin{equation}
			\forall r, s, t \in R: (r + s) \cdot t = r \cdot t + s \cdot t
		\end{equation}
	\end{enumerate}
\end{definition}

\begin{definition}[Commutative Ring]
	A commutative ring $\{ R, +, \cdot \}$ is a ring where $\{ R, \cdot \}$ is commutative.
\end{definition}

\begin{definition}[Field]
	A field $\{ \mathbb{F}, +, \cdot \}$ is a ring where $\{ \mathbb{F} \setminus 0, \cdot \}$ is a multiplicative group.
\end{definition}

\section{Module-Like Structures}

Modules generalize vector spaces.

\begin{definition}[Left Module]
	A left module $A$ over a ring $R$ is a system of sets and operations $\{ A, R, \oplus, +, \cdot, * \}$ which satisfy the following properties.
	\begin{enumerate}
		\item $\{ A, \oplus \}$ is an additive Abelian group.
		\item $\{ R, +, \cdot \}$ is a ring.
		\item Scalar multiplication $*: R \times A \rightarrow A$ is a left action of the multiplicative monoid $\{ R, \cdot \}$ of the ring on $A$ which satisfies the following properties.
		\begin{enumerate}
			\item Distributivity over $A$'s addition
			\begin{equation}
				\forall r \in R, a, b \in A: r * (a \oplus b) = r * a \oplus r * b
			\end{equation}
			\item Distributivity over $R$'s addition into $A$'s addition
			\begin{equation}
				\forall r, s \in R, a \in A: (r + s) * a = r * a \oplus s * a
			\end{equation}
		\end{enumerate}
	\end{enumerate}
\end{definition}

\begin{definition}[Left Vector Space]
	A left vector space $V$ over a field $\mathbb{F}$ is a left module with sets and operations $\{V, \mathbb{F}, \oplus, +, \cdot, *\}$ where the ring is a field. Scalar multiplication is then a left action of the multiplicative group $\{ \mathbb{F} \setminus 0, \cdot \}$.
\end{definition}

Further, scalar multiplication can be used to relate the inverse and identity elements of the field and vector space.

\begin{flalign}
	&& \forall \vb{v} \in V: 0 * \vb{v} &= \vb{0} & \text{Proof \ref{proof:smult_zero}} \label{equation:smult_zero} \\
	&& \forall \vb{v} \in V: (-1) * \vb{v} &= -\vb{v} & \text{Proof \ref{proof:smult_inverse}} \label{equation:smult_inverse}
\end{flalign}

In practice, the notation used is looser. Vector and scalar addition share the same symbol. Field and scalar multiplication symbols are usually dropped. The operations can be inferred from context.

\section{Derived Vector Spaces}

The set of functions $V^{X}$ mapping from an arbitrary set $X$ to a vector space $V$ over $\mathbb{F}$ can be turned into a vector space over $\mathbb{F}$ with the following operations.

\begin{flalign}
	&& (f + g)(x) &= f(x) + g(x) & \\
	&& (\alpha * f) &= \alpha * f(x) &
\end{flalign}

\begin{definition}[Linear Subspace]
	A linear subspace $U \subset V$ of a vector space $V$ over $\mathbb{F}$ is a subset which is also a vector space using $V$'s operations. This means that it is closed with respect to vector addition and scalar multiplication.
\end{definition}

Affine spaces generalize spaces with no inherent reference point. For example consider Euclidean space or time. There is no inherent coordinate system or time axis to the universe. Assigning a coordinate origin, calendar system, or time zone is arbitrary. However, once this is done, the affine space can be represented with a vector space.

\begin{definition}[Affine Space]
	An affine space $\{ X, V, +, \oplus \}$ is a set $X$, vector space $V$ over a field $\mathbb{F}$, and a right action $\oplus: X \times V \rightarrow X$ of the additive group $\{ V, + \}$ which satisfies the following properties.
	\begin{enumerate}
		\item The action is
		\begin{enumerate}
			\item Free
			\begin{equation}
				\forall \vb{v}\in V: x \oplus \vb{v} = x \implies \vb{v} = \vb{0}
			\end{equation}
			\item Transitive
			\begin{equation}
				\forall x, y \in X: \exists \vb{v} \in V: x \oplus \vb{v} = y
			\end{equation}
		\end{enumerate}
	\end{enumerate}
\end{definition}

Elements in $X$ are called points and vectors in $V$ are called translations or free vectors.

The definition above also implies the following properties.

\begin{flalign}
	&& \forall x, y \in X: \exists! \vb{v} \in V: x \oplus \vb{v} &= y & \text{Bijection  for fixed $x$. Proof \ref{proof:aff_bij_x}} \label{equation:aff_bij_x} \\
	&& \forall \vb{v} \in V, y \in X: \exists! x \in X: x \oplus \vb{v} &= y & \text{Bijection for fixed $\vb{v}$. Proof \ref{proof:aff_bij_v}} \label{equation:aff_bij_v}
\end{flalign}

Since these operations are invertible, both point-vector $\ominus: X \times V \rightarrow X$ and point-point $\dot{-}: X \times X \rightarrow V$ subtraction operations are well-defined (Proof \ref{proof:aff_sub}).

\begin{equation}
	a \oplus (b \dot{-} a) = b
\end{equation}

This operation also satisfies the following properties, called Weyl's axioms.

\begin{flalign}
	&& \forall x \in X, \vb{v} \in V: \exists! y \in x: y \dot{-} x &= \vb{v} & \text{Proof \ref{proof:aff_w1}} \label{equation:aff_w1} \\
	&& \forall x, y, z \in X: (c \dot{-} b) + (b \dot{-} a) &= c \dot{-} a & \text{Proof \ref{proof:aff_w2}} \label{equation:aff_w2}
\end{flalign}

Furthermore, the parallelogram property is satisfied.

\begin{equation}
	\forall w, x, y, x \in X: b \dot{-} a = d \dot{-} c \Longleftrightarrow c \dot{-} a = d \dot{-} b
\end{equation}

As with vector spaces, distinguishing all of these operations with special symbols is done for clarity. In practice, simple addition and subtraction are used, and such notation will be used moving forwards.

Finally, any vector space is an affine space over itself (Proof \ref{proof:vec_self_aff}). This means that any element of $V$ can be considered as a point or a vector. Importantly, this formalizes that notion of a reference point being arbitrary. Given a reference point $o \in A$, an isomorphism $\iota_{o}: X \leftrightarrow V$ can be constructed between $A$ and $V$.

\begin{equation}
	\iota_{o}(x) = x \ominus o
\end{equation}

\begin{definition}[Affine Subspace]
	An affine subspace $W \subseteq X$ of an affine space $\{ X, V, +, \oplus \}$ is a set $W$ for which there exists some $w \in W$ such that generated vectors $\overrightarrow{W} = \{ x \ominus w: x \in W \}$ is a linear subspace of V.
\end{definition}

It folllows that the choice of $w$ is arbitrary, and all subspaces can be generated from linear subspaces $U \subseteq V$ as follows.

\begin{equation}
	x + U := \{ x + \vb{u}: \vb{u} \in U \}
\end{equation}

The linear subspace $U$ can be considered the affine subspace's direction, and two subspaces that share the same or subsidiary directions $T \subseteq U$ are considered parallel. The above properties imply Playfair's axiom: for any point $x \in X$ and direction $U \subseteq V$ there is a unique affine subspace of direction $U$ (Proof \ref{proof:playfair}).

\begin{definition}[Affine Map]
	An affine map $f: A \rightarrow B$ between two affine spaces $A$ and $B$ is one which satisfies the following property.
	
	\begin{equation}
		\forall a, b, c, d \in A: a \dot{-} b = c \dot{-} d \implies f(a) \dot{-} f(b) = f(c) \dot{-} f(d)
	\end{equation}
\end{definition}

This implies that the following map between the associated vector spaces $\overrightarrow{f}: \overrightarrow{A} \rightarrow \overrightarrow{B}$ is well-defined linear map (Proof \ref{proof:vec_aff_welldef}).

\begin{equation}
	\overrightarrow{f}(b \dot{-} a) = f(b) \dot{-} f(a)
\end{equation}

The flip side of this is that one can start with arbitrary points $a \in A$ and $b \in B$ and a linear map $\overrightarrow{f}$ and generate $f$.

\begin{equation}
	f(x) = b \oplus \overrightarrow{f}(x \dot{-} a)
\end{equation}

\section{Basis}

\begin{definition}[Linear Combination]
	A linear combination of vectors $\{ \vb{v}_{j} \}_{j = 1}^{k}$ is a finite sum of scalar multiples of themselves.
	
	\begin{equation}
		\sum_{j = 1}^{k} \alpha_{j} \vb{v}_{k}
	\end{equation}
\end{definition}

\begin{definition}[Linear Dependence]
	A vector $\vb{v}$ is linearly dependent on a set of vectors $\{ \vb{v}_{j} \}_{j = 1}^{k}$ if there is a linear combination of the set which equals $v$.
	
	\begin{equation}
		\exists \{ \alpha_{j} \}_{j = 1}^{k} \subseteq \mathbb{F}: \vb{v} = \sum_{j = 1}^{k} \alpha_{j} \vb{v}_{k}
	\end{equation}
\end{definition}

\begin{definition}[Linear Independence]
	A set of vectors $\{ \vb{v}_{j} \}_{j = 1}^{k}$ is linearly independent if none of the vectors are linearly dependent on the others. Otherwise, the set is linearly dependent. An infinite set is linearly independent if every finite subset of it is linearly independent.
\end{definition}

\begin{definition}[Hamel Basis]
	A Hamel basis is a maximal linearly independent subset of a vector space. This means that it is not a proper subset of another linearly independent subset.
\end{definition}

If it exists, a Hamel basis spans its vector space and generates unique finite linear combinations for every vector in the space.

\begin{definition}[Span]
	A subset $B$ spans a vector space $V$ if every vector in $V$ is linearly dependent on a finite subset of $B$.
\end{definition}

\begin{definition}[Dimension]
	The dimension of a vector space is the cardinality of its basis.
	
	\begin{equation}
		\dim{V} = \# B
	\end{equation}
\end{definition}

\subsection{Topological Functions}

\begin{definition}[Metric]
	A metric $d: X \times X \rightarrow [0, \infty)$ is a function that generalizes distance between elements of the set $X$ by satisfying the following properties.
	\begin{enumerate}
		\item Zero property
		\begin{equation}
			\forall x, y \in X: d(x, y) = 0 \iff x = y
		\end{equation}
		\item Symmetry
		\begin{equation}
			\forall x, y \in X: d(x, y) = d(y, x)
		\end{equation}
		\item Triangle inequality
		\begin{equation}
			\forall x, y, z \in X: d(x, y) \leq d(x, z) + d(y, z)
		\end{equation}
	\end{enumerate}
\end{definition}

A set with an associated metric is called a metric space.

\begin{definition}[Norm]
	A norm $\norm{\cdot}: V \rightarrow [0, \infty)$ is a function that generalizes magnitude of vectors in a vector space $X$ over a real or complex field $\mathbb{F} \subseteq \mathbb{C}$ by satisfying the following properties.
	\begin{enumerate}
		\item Positive-definiteness
		\begin{equation}
			\forall \vb{v} \in V: \norm{\vb{v}} = 0 \iff \vb{v} = \vb{0}
		\end{equation}
		\item Absolute homogeneity, i.e. scaling property
		\begin{equation}
			\forall \alpha \in \mathbb{F}, \vb{v} \in V: \norm{\alpha \vb{v}} = \abs{\alpha} \norm{\vb{v}}
		\end{equation}
		\item Triangle inequality
		\begin{equation}
			\forall \vb{u}, \vb{v} \in V: \norm{\vb{u} + \vb{v}} \leq \norm{\vb{u}} + \norm{\vb{v}}
		\end{equation}
	\end{enumerate}
\end{definition}

A vector space with a norm is called a normed linear space. It is also a metric space using the induced metric.

\begin{flalign}
	&& d(\vb{u}, \vb{v}) &= \norm{\vb{u} - \vb{v}} & \text{Proof \ref{proof:nls_is_ms}} \label{equation:nls_is_ms}
\end{flalign}

All norms are equivalent in finite-dimensional spaces. See Proof \ref{proof:norm_equivalence}.

\begin{definition}
	An inner product $(\cdot, \cdot): V \times V \rightarrow \mathbb{F}$ is a function that generalizes the similarity of vectors in a vector space $X$ over a real or complex field $\mathbb{F} \subseteq \mathbb{C}$ by satisfying the following properties.
	\begin{enumerate}
		\item Positive-definiteness
		\begin{equation}
			\forall \vb{x} \in X: (\vb{x}, \vb{x}) = 0 \Longleftrightarrow \vb{X} = \vb{0}
		\end{equation}
		\item Linearity in the first argument
		\begin{equation}
			\forall \alpha_{1}, \alpha_{2} \in \mathbb{F}, \vb{x}_{1}, \vb{x}_{2}, \vb{y} \in X: (\alpha_{1} \vb{x}_{1} + \alpha_{2} \vb{x}_{2}, \vb{y}) = \alpha_{1} (\vb{x}_{1}, \vb{y}) + \alpha_{2} (\vb{x}_{2}, \vb{y})
		\end{equation}
		\item Conjugate symmetry
		\begin{equation}
			\forall \vb{x}, \vb{y} \in X: (\vb{x}, \vb{y}) = \overline{(\vb{y}, \vb{x})}
		\end{equation}
	\end{enumerate}
\end{definition}

These conditions imply anti-linearity in the second argument, or linearity for functions over a real field.

\begin{flalign}
	&& \forall \beta_{1}, \beta_{2} \in \mathbb{F}, \vb{x}, \vb{y}_{1}, \vb{y}_{2} \in X: (\vb{x}, \beta_{1} \vb{y}_{1} + \beta_{2} \vb{y}_{2}) &= \overline{\beta_{1}} (\vb{x}, \vb{y}_{1}) + \overline{\beta_{2}} (\vb{x}, \vb{y}_{2}) & \text{Proof \ref{proof:ip_sesqui}} \label{equation:ip_sesqui}
\end{flalign}

For any inner product, the Cauchy-Schwarz inequality applies.

\begin{flalign}
	&& (\vb{x}, \vb{y}) &\leq \sqrt{(\vb{x}, \vb{x}) (\vb{y}, \vb{y})} & \text{Proof \ref{proof:cs_ineq}} \label{equation:cs_ineq}
\end{flalign}

A vector space with an inner product is called an inner product space. It is also a normed linear space and a metric space using the induced norm.

\begin{flalign}
	&& \norm{\vb{x}} &= \sqrt{(\vb{x}, \vb{x})} & \text{Proof \ref{proof:ips_is_nls}} \label{equation:ips_is_nls}
\end{flalign}

An inner product also allows for an abstract notion of angle, inspired by geometry.

\begin{equation}
	\cos(\theta) = \frac{(\vb{x}, \vb{y})}{\norm{\vb{x}} \norm{\vb{y}}}
\end{equation}

\section{Dual Space}

\begin{definition}[Algebraic Dual Space]
content...
\end{definition}

\begin{definition}[Topological Dual Space]
content...
\end{definition}

\section{Riesz Representation}

Riesz representation and dealing with complex vector spaces, dual basis cobasis link, also compare components

\subsection{Cobasis}

\section{Tensor Product Spaces}

over underline otimes, 4th order transpose?

\subsection{Action of a Tensor Product}

\subsection{Einstein Index Notation}